\subsection{Policy Implications for Digital Inclusion}

Our findings have direct implications for broadband policy in Europe and neighboring regions. The strong pre-COVID price elasticity in EaP countries suggests that affordability-focused policies---price regulation, operator subsidies, or consumer vouchers---can effectively increase broadband adoption. A 10\% price reduction would increase subscriptions by approximately $\EaPTenPctChg$\%, translating to nearly \EaPAddSubsTotal{} new users across the EaP region.

However, the shift toward inelastic demand during COVID-19 fundamentally alters the policy toolkit. When price sensitivity approaches zero, further price reductions yield minimal adoption gains. Instead, policies must address other barriers including infrastructure gaps (particularly in rural areas), digital literacy deficits, and content relevance \citep{guermazi2021broadband}. The transition from price-sensitive to price-insensitive demand marks broadband's evolution from discretionary service to essential utility.

For EU countries, the consistently low elasticity even pre-COVID suggests price is not the primary adoption barrier. With subscription rates exceeding \DescSubsEU{} per 100 inhabitants (Table~\ref{tab:descriptives}), most households with willingness to pay already have broadband. Remaining gaps likely reflect supply constraints (rural coverage) or demographic factors (elderly populations). EU policy should focus on universal service obligations ensuring infrastructure availability rather than price interventions.

The EaP--EU differential carries an important temporal dimension. Prior to 2020, the gap in elasticities ($\EaPBaseBr$ vs.\ $\EUBaseBr$) implied that affordability-focused interventions were far more potent in EaP than in EU markets, justifying differentiated policy designs within the EU4Digital framework. However, our central finding is that this differential has collapsed: by 2020--2024 both regions exhibit near-zero elasticity and the case for price-based instruments has weakened simultaneously across the entire sample. The policy implication is not simply that EaP and EU need different instruments, but that \emph{both} require a shift away from affordability programs toward infrastructure deployment and digital skills investment as broadband matures into an essential utility.

\subsection{Time-Varying Elasticity and Demand Evolution}

A central finding is that broadband demand elasticity is not constant but evolves systematically over time. For EU countries (N$=$27), the year-by-year analysis reveals a statistically precise declining trend starting well before COVID-19, with elasticity falling from $\EUYbyYRefElast$ in 2015 to near zero by the late 2010s. For EaP countries, year-by-year estimates cannot be interpreted at the annual frequency: with only five to six EaP countries, each triple interaction term is identified from too few cross-sectional units and is confounded with EaP-bloc-specific shocks (geopolitical events, investment cycles) that cannot be separated from the price response. Instead, the pooled placebo test (Online Appendix Table~A4) confirms that EaP elasticity was also declining before COVID-19 ($+\PlaceboEaP$, p$=$\PlaceboEaPPval), as was EU elasticity ($+\PlaceboEU$, p$=$\PlaceboEUPval). This gradual transformation in both regions suggests demand elasticity is a function of market maturity, technological change, and societal integration.

Three mechanisms likely drive this evolution. First, \textbf{network effects} become stronger as broadband penetration increases. The value of connectivity rises with the number of connected users, making broadband more essential regardless of price \citep{katz1985network}. Second, \textbf{habit formation} makes broadband increasingly indispensable as consumers integrate it into daily routines \citep{becker1988theory}. Third, \textbf{complementary innovations}---smartphones, streaming services, cloud computing---increase broadband's value proposition independent of price.

This time-varying elasticity has important methodological implications. Static elasticity estimates from early-period data may not apply to current markets. Panel methods allowing for time variation, as employed here, are essential for capturing evolving demand patterns.

From a forecasting perspective, our results suggest broadband adoption will continue even without further price declines. The near-zero COVID-era elasticity implies that infrastructure availability and service quality---not affordability---now determine adoption. Countries planning broadband expansion should prioritize network deployment over price regulation as markets mature.

\subsection{Price Measurement and Affordability Metrics}

The robustness across price measures (Online Appendix Table~A3) provides valuable methodological insights. All three price definitions---income-relative (GNI\%), PPP-adjusted, and nominal USD---yield negative elasticity estimates, confirming that broadband demand is price-sensitive regardless of measurement choice. However, the magnitudes differ substantially: income-relative prices yield the largest EaP elasticities ($\EaPGNIFCB$), followed by nominal USD ($\EaPUSDBCB$), and PPP-adjusted ($\EaPPPPFCB$).

However, \textit{income-relative prices} emerge as the most appropriate affordability metric. The ITU basket is designed to track the minimum viable entry-level broadband offering in each period---1\,GB in the early years, 5\,GB once that became the prevailing threshold---so the GNI\% measure is literally the fraction of monthly income a household must spend to cross the adoption threshold. This directly matches the object of interest in a discrete adoption model: not how expensive broadband is in absolute terms, but whether it is affordable enough to subscribe at all \citep{deaton1980almost}. A \$30/month connection is inexpensive for a household earning \$3{,}000/month but prohibitive for one earning \$200/month; GNI\% captures this heterogeneity while nominal USD does not.

This framing also clarifies the PPP failure. PPP deflators are constructed from consumer expenditure baskets dominated by food, 
housing, and tradable goods---categories where quality-adjusted prices are stable. 
Technology services, by contrast, exhibit rapid quality-adjusted price declines \citep{schreyer2002oecd}. 
Applying PPP to broadband prices therefore introduces a systematic deflation error, 
understating the purchasing-power cost of connectivity in lower-income countries and 
reducing the precision of cross-country affordability comparisons. 
Income-relative pricing sidesteps this entirely by using income---rather than a mismeasured price index---as the denominator.

Empirically, income-relative prices show the strongest relationship with subscriptions (highest $R^2$ values across specifications), 
as demonstrated in Figure~\ref{fig:price_measurement}. This validates the ITU's practice of expressing 
affordability targets as percentages of GNI per capita \citep{itu2024data}. Policymakers tracking broadband affordability 
should prioritize income-relative metrics over both nominal and PPP-adjusted prices.

As noted in Section~\ref{sec:data}, the ITU basket definition changed in 2018; because the 
transition was simultaneous across all countries, year fixed effects absorb the common level 
shift and the redefinition does not confound within-country elasticity estimates. 
Actual consumer prices vary by speed tier, data caps, and bundling arrangements beyond what standardized baskets capture. 
To the extent measurement error exists, classical error would attenuate estimates toward zero, 
making our significant negative elasticities conservative. 
Future research could refine price measurement using micro-level tariff data, 
though cross-country comparability would be challenging.

\begin{figure}[t]
\centering
\includegraphics[width=0.75\textwidth]{fig3_price_measurement.pdf}
\caption{Effect of price measurement on inference (pre-COVID, 2010--2019). (a) Mean elasticity by price measure. (b) Percentage of specifications significant at p$<$0.05. GNI\% yields 100\% significant EaP results vs.\ \PctPPPSigEaP\% for PPP (at p$<$0.05), validating income-relative pricing as the appropriate affordability metric.}
\label{fig:price_measurement}
\end{figure}

\subsection{Regional Heterogeneity and Development Gradients}

The large EaP-EU elasticity differential reflects broader development gradients. EaP countries have lower GDP per capita (\$4,500--\$13,000 PPP) compared to EU countries (\$25,000--\$60,000 PPP), translating directly to affordability constraints and higher price sensitivity.

Beyond income, institutional and infrastructure differences matter. EaP countries exhibit lower regulatory quality (Table~\ref{tab:descriptives}), potentially reducing competitive pressure and raising prices. Weaker governance may also hamper policy implementation of broadband programs. Additionally, lower internet server density and R\&D intensity suggest less developed digital ecosystems, reducing complementary services that drive demand.

These structural differences imply that development strategies successful in high-income EU markets may not transfer directly to middle-income EaP contexts. Targeted approaches addressing affordability constraints and institutional gaps are essential.

Interestingly, the convergence toward near-zero elasticity during COVID occurred in both regions, suggesting the pandemic's impacts transcended development levels. This may reflect the universal need for connectivity during social restrictions, overriding typical income-based differences in demand. Crucially, this convergence constitutes the first econometric evidence of achieved demand-side digital harmonisation between EaP and EU markets: the policy objective articulated under the EU4Digital initiative---that EaP consumers come to treat broadband as they do in the EU, as a non-discretionary essential service---is not a distant aspiration but an observable outcome already realised in the data.

\subsection{COVID-19 Shock vs. Secular Trends}

A crucial interpretive question is whether observed changes reflect a COVID-19 shock or
continuation of pre-existing trends. The placebo test (Online Appendix Table~A4) provides
suggestive evidence for the latter: both EU and EaP
placebo interactions are positive and marginally significant, indicating that elasticity
was already declining during 2015--2019.

The period-split results (Online Appendix Table~B1) sharpen this interpretation by
decomposing the post-2019 period into two distinct phases. During the \emph{acute} COVID
phase (2020--2021), lockdowns and remote-work mandates plausibly converted broadband from
a convenience to a necessity, suppressing price sensitivity. During the \emph{post-acute}
phase (2022--2024), pandemic restrictions eased and work/education patterns largely normalized
in most sample countries. Yet near-zero elasticity persisted. This persistence is
inconsistent with a pure pandemic-shock interpretation---if the shock were the mechanism,
one would expect some reversion toward pre-COVID elasticity as restrictions lifted.

Instead, zero post-acute elasticity supports the \emph{secular transformation} hypothesis:
broadband has permanently crossed the threshold from discretionary to essential service.
High-speed connectivity is now deeply embedded
in work, education, healthcare, and social life. Even at elevated prices, households treat
disconnection as prohibitively costly---analogous to electricity or water utilities. COVID-19
did not cause this transition; it accelerated the final stage of a process underway since
roughly 2015.

For policy, this distinction has permanent implications. If zero elasticity is durable,
price subsidies lose their effectiveness. The binding
constraints are now infrastructure availability (especially in rural areas),
digital literacy, and service quality. EU policy should focus
on universal coverage; EaP policy should pair infrastructure investment with
institutional reforms to foster competition.

\subsection{Limitations}

Several limitations qualify our findings. Country-level data mask within-country heterogeneity (urban--rural gaps, income distribution); household-level data would enable richer analysis but are unavailable cross-country over 15 years. Our focus on fixed broadband offers the longest consistent time-series, though mobile substitution is a concern: adding log mobile subscriptions as a control leaves elasticity unchanged (EU: $\EUBaseBr$; EaP: $\EaPBaseBr$), confirming the estimates are not confounded by mobile expansion. ITU price baskets do not capture quality heterogeneity (speed, reliability), and lack of quality-adjusted prices is standard in cross-country telecommunications research \citep{greenstein2016measuring}. Endogeneity is addressed directly via 2SLS (Online Appendix Table~B2; first-stage $F = \IVFirstF$ pre-COVID): 2SLS estimates are virtually identical to OLS, confirming simultaneity bias is negligible. Finally, generalisability beyond Europe and the Eastern Partnership requires caution, as elasticity likely varies with income and institutions.

Three alternative explanations for near-zero COVID-era elasticity---operator capacity constraints, within-year timing mismatch in annual data, and government price interventions \citep{oecd2021broadband}---cannot be fully ruled out but are inconsistent with the evidence of gradual pre-COVID decline beginning in 2015. Robustness across specifications and price measures (Online Appendix Table~A3) provides confidence in the core findings.
